\documentclass[10pt,a4paper]{amsart}
\usepackage[margin=19mm]{geometry}
\usepackage{amsmath,amssymb,mathtools}
\usepackage[scr=rsfs,cal=euler]{mathalfa}
\usepackage{xfrac}
\usepackage[unicode,hidelinks,bookmarks=false]{hyperref}
\newcommand{\ie}{i.e.\ }
\newcommand{\cf}{cf.\ }
\newcommand{\resp}{respectively\ }
\newcommand{\Subsection}{\subsection}
\providecommand{\nobreakdash}{\nobreak\hbox{-}}
\setlength{\emergencystretch}{2em}
\multlinegap0pt
\usepackage{amssymb}
\usepackage{xcolor}
\newtheorem{lemma}{Lemma}
\newcommand{\E}{\mathbb E}
\title{Sharp Diameter Bounds for Nonnegative\\Cyclotomic Multiples}
\author{Hu Tan, Ying Zhang}
\date{Source: arXiv:2609.07646v1 (CC BY 4.0)}
\begin{document}
\maketitle

\noindent\emph{Source and licence.} Sharp Diameter Bounds for Nonnegative\\Cyclotomic Multiples. Version arXiv:2609.07646v1 (CC BY 4.0), distributed by arXiv under the Creative Commons Attribution 4.0 International licence, http://creativecommons.org/licenses/by/4.0/, as recorded in the arXiv licence metadata for this exact version and shown on the abstract page https://arxiv.org/abs/2609.07646v1. Full licence notice: \texttt{licenses/arxiv-2609.07646-CC-BY-4.0.txt}.\par\smallskip

\noindent\textit{Editorial scope.} Counted unit: the article's lemma lem:arc (source0.tex lines 216--231). The article's complete original proof of it is delivered with it (source0.tex lines 233--292, one \texttt{proof} environment of 60 lines). No mathematics is written, repaired or re-derived: the statement and the proof are the article's own lines, in the article's own notation, and no retained line is substituted or re-flowed.  No further source-local context is needed: every cross-reference of every retained line resolves inside the two delivered spans.  Nothing was omitted from any retained span, so the package carries no omission record.  No retained line contains a citation, so the wrapper carries no bibliography and no bibliography entry.  No retained line contains a figure environment, an image-inclusion command or drawing code, so no image is delivered and the package declares no asset dependency.  Editorial formatting only, in addition: an AMS \texttt{amsart} preamble that restates the article's own declarations and macros used by the retained lines, namely the environments \texttt{lemma} on separate counters, together with the standard packages those lines need.  The retained environments print in this document as Lemma 1 in order of appearance, whereas the article prints the same items with the numbering of its own full document.  The complete native source of the article as distributed in the arXiv e-print for this exact version is bundled unchanged as \texttt{sources/209646/source0.tex}.
\begin{lemma}[Arc bound]\label{lem:arc}
Let $r\ge2$ be an integer. Suppose that a nonzero finite positive measure
$\mu$ is supported on $[-\alpha,\alpha]$, where $0\le\alpha<\pi$, and
\begin{equation}\label{eq:moments}
 \widehat\mu(k)=0\qquad(1\le k\le r-1).
\end{equation}
Then
\begin{equation}\label{eq:arc-bound}
 \alpha\ge\pi-\frac\pi r.
\end{equation}
If equality holds, $\mu$ assigns equal positive masses to the $r$ points
\begin{equation}\label{eq:nodes}
 -\alpha+\frac{2\pi j}{r},\qquad 0\le j\le r-1,
\end{equation}
and has no other support.
\end{lemma}

\begin{proof}
Normalize $\mu$ to be a probability measure, and write
$\E f=\int f\,d\mu$. Set
\begin{equation}\label{eq:sine-vector}
 t=\frac\pi r,\qquad
 v_j=(-1)^j\sin((j+1)t)\quad(0\le j\le r-2),
\end{equation}
and define
\begin{equation}\label{eq:Q}
 Q(z)=\sum_{j=0}^{r-2}v_jz^j,\qquad
 S=\sum_{j=0}^{r-2}v_j^2=\frac r2.
\end{equation}
The last equality is the elementary sum
$\sum_{j=1}^{r-1}\sin^2(j\pi/r)=r/2$.
The vanished moments give
\begin{align}
 \E|Q(e^{i\theta})|^2&=S,\label{eq:norm}\\
 \E\!\left[\cos\theta\,|Q(e^{i\theta})|^2\right]
 &=\sum_{j=0}^{r-3}v_jv_{j+1}
 =-\cos(t)S.\label{eq:cosine}
\end{align}
Indeed, the trigonometric polynomials in these two expectations have
degrees at most $r-2$ and $r-1$, respectively, so only their constant
Fourier coefficients remain. The final identity follows by multiplying
\[
 v_{j-1}+v_{j+1}=-2\cos(t)v_j\quad(0\le j\le r-2),
 \qquad v_{-1}=v_{r-1}=0,
\]
by $v_j$ and summing. This also covers $r=2$, when the adjacent sum
is empty.

On the support of $\mu$, $\cos\theta\ge\cos\alpha$. Hence
\[
 -\cos(t)S\ge\cos\alpha\,S.
\]
Since $S>0$ and cosine is decreasing on $[0,\pi]$, this proves
\eqref{eq:arc-bound}.

If $\alpha=\pi-t$, equality in this integral inequality forces the
support into the zero set of
\begin{equation}\label{eq:H}
 (\cos\theta+\cos t)|Q(e^{i\theta})|^2
 \quad\text{on }[-\alpha,\alpha].
\end{equation}
The recurrence above gives the polynomial identity
\begin{equation}\label{eq:Qidentity}
 (1+2\cos(t)z+z^2)Q(z)=\sin(t)\bigl(1+(-z)^r\bigr).
\end{equation}
The polynomial $1+(-z)^r$ has the $r$ distinct roots
$e^{i(-\alpha+2\pi j/r)}$. The two roots of the quadratic on the left
are $e^{\pm i\alpha}$, and the remaining $r-2$ roots are the roots of
$Q$. Thus the zeros in \eqref{eq:H} are exactly \eqref{eq:nodes}.
If $w_j$ denotes the mass at the $j$th node, allowing initially $w_j=0$,
then
\[
 \sum_{j=0}^{r-1}w_je^{2\pi ikj/r}=0\quad(1\le k<r),
 \qquad \sum_{j=0}^{r-1}w_j=1.
\]
Discrete Fourier inversion gives $w_j=1/r$ for all $j$.
\end{proof}

\end{document}
